\section*{Bab \ref{ch:b3:representations} --- Representasi
Grup Berhingga}

\begin{solution}{exo:b3:representations:1}
(a) Grup $\Z/n\Z$ abelian: jadi semua \omterm{def:b3:representations:rep}{representasi tak tereduksinya} berderajat $1$
(\cref{prop:b3:representations:onedim}), yakni berupa morfisma
$\chi \colon \Z/n\Z \to \C^\times$, yang ditentukan oleh $\chi(\bar 1)
= \omega$ dengan $\omega^n = 1$: jadi ada $n$ \omterm{def:b3:representations:character}{karakter} $\chi_k(\bar
m) = \eu^{2\iu\pi km/n}$. Untuk $n = 4$ (kelasnya $=$ unsurnya
$\bar0,\bar1,\bar2,\bar3$):
\[
\begin{array}{c|cccc}
& \bar0 & \bar1 & \bar2 & \bar3\\
\hline
\chi_0 & 1 & 1 & 1 & 1\\
\chi_1 & 1 & \iu & -1 & -\iu\\
\chi_2 & 1 & -1 & 1 & -1\\
\chi_3 & 1 & -\iu & -1 & \iu
\end{array}
\]
(b) Barisnya: $\langle\chi_k,\chi_l\rangle = \frac14\sum_m
\eu^{2\iu\pi(l-k)m/4} = \delta_{kl}$ (lewat jumlah geometri); dan kolomnya
serupa. Matriks $(\eu^{2\iu\pi km/n})_{k,m}$ adalah matriks
\emph{transformasi Fourier diskret} --- yakni penyaring
akar satuan yang sama seperti pada bab fungsi pembangkit di
jilid Tahun ke-2; jadi \omterm{thm:b3:representations:orthogonality}{ortogonalitas karakter}
menyamaratakan rumus inversi transformasi Fourier diskret itu.
\end{solution}

\begin{solution}{exo:b3:representations:2}
(a) Matriks $\rho(g)$ dalam basis $(e_x)$ berupa
matriks permutasi, yang tracenya sama dengan banyaknya $x$ dengan $g\cdot x =
x$. Maka
\[
\langle\mathbf 1, \chi\rangle = \frac1{\abs G}\sum_g
\abs{\operatorname{Fix}(g)} = \#\{\text{orbit}\}
\]
menurut lema pencacahan Burnside (\cref{exo:b3:groups:5}) ---
setara dengan itu, perhitungan ini memberikan multiplisitas \omterm{def:b3:representations:rep}{representasi}
trivialnya, yang ruang isotipiknya adalah ruang vektor
invarian-$G$, berdimensi sama dengan banyaknya \omterm{def:b3:groups:action}{orbit} (satu
fungsi indikator per \omterm{def:b3:groups:action}{orbit}).

(b) Karena $\operatorname{Fix}_{X\times X}(g) =
\operatorname{Fix}_X(g)^2$, bagian (a) yang diterapkan pada $X \times X$
memberikan $\langle\chi,\chi\rangle = \frac1{\abs
G}\sum\abs{\operatorname{Fix}(g)}^2 = \#\{\text{orbit pada }
X\times X\}$ (sebab $\chi$ real). Dengan menulis $\chi = \mathbf 1 +
\psi$: $\langle\mathbf1,\chi\rangle = 1$ (karena transitif), sehingga
$\langle\psi,\psi\rangle = \langle\chi,\chi\rangle - 1$. Aksi
pada $X\times X$ mempunyai diagonalnya sebagai satu \omterm{def:b3:groups:action}{orbit}; dan tepat
ada satu \omterm{def:b3:groups:action}{orbit} lain jika dan hanya jika $G$ transitif pada pasangan yang berbeda:
jadi $\langle\psi,\psi\rangle = 1$ jika dan hanya jika aksinya $2$-transitif
(\cref{cor:b3:representations:multiplicity}).

(c) Grup $S_n$ bersifat $2$-transitif pada $\{1,\dots,n\}$ (sebab ia mengirim sembarang
pasangan berbeda ke mana pun): jadi $\psi = \chi_{\mathrm{std}}$
\omterm{def:b3:representations:rep}{tak tereduksi}.
\end{solution}

\begin{solution}{exo:b3:representations:3}
$S_3$: ada tiga kelas dengan $\sum n_i^2 = 6 = 1 + 1 + 4$; kedua
\omterm{def:b3:representations:character}{karakter} berderajat $1$ adalah $\mathbf 1, \varepsilon$ (sebab $[S_3 :
A_3] = 2$); sedangkan baris ketiganya $(2, a, b)$ menyusul dari ortogonalitas
kolom dengan kolom-$e$: $1 - 1 + 2a = 0$ dan $1 + 1 +
2b = 0$: jadi $a = 0$, $b = -1$ --- persis tabel
\cref{ex:b3:representations:s3}.

Pemeriksaan $S_4$ (barisnya): $\langle\chi_{\mathrm{std}},
\varepsilon\chi_{\mathrm{std}}\rangle = \frac1{24}(9 - 6 + 3 + 0
- 6) = 0$; dan $\langle\chi_2,\chi_2\rangle = \frac1{24}(4 + 0 + 12
+ 8 + 0) = 1$. Kolomnya: $e$ terhadap $(1\,2)$: $1 - 1 + 0 + 3 - 3
= 0$; lalu $(1\,2)$ terhadap dirinya sendiri: $1 + 1 + 0 + 1 + 1 = 4 =
\abs{Z_{S_4}((1\,2))} = 24/6$.

\omterm{def:b3:representations:character}{Karakter} permutasi pada $4$ titik: $(4, 2, 0, 1, 0) = \mathbf
1 + \chi_{\mathrm{std}}$ (lewat cacahan titik tetap; kurangkan baris
teratasnya). Untuk $\chi_{\mathrm{std}}^2 = (9, 1, 1, 0, 1)$:
\[
\langle\mathbf1,\cdot\rangle = \tfrac{9 + 6 + 3 + 0 + 6}{24} =
1,\quad
\langle\varepsilon,\cdot\rangle = 0,\quad
\langle\chi_2,\cdot\rangle = \tfrac{18 + 6}{24} = 1,\quad
\langle\chi_{\mathrm{std}},\cdot\rangle = 1,\quad
\langle\varepsilon\chi_{\mathrm{std}},\cdot\rangle = 1:
\]
$\chi_{\mathrm{std}}^2 = \mathbf 1 + \chi_2 +
\chi_{\mathrm{std}} + \varepsilon\chi_{\mathrm{std}}$
(dimensinya: $9 = 1 + 2 + 3 + 3$).
\end{solution}

\begin{solution}{exo:b3:representations:4}
Kedua grup itu mempunyai lima kelas dan pola derajat $(1,1,1,1,2)$
(empat \omterm{def:b3:representations:character}{karakter} berderajat $1$ dari abelianisasinya yang $\cong (\Z/2\Z)^2$,
lalu $\sum n_i^2 = 8$). Dengan mengurutkan kelasnya $e$, $z$ (yaitu
involusi pusatnya: $r^2$, atau $-1$), dan ketiga
kelas beranggota dua:
\[
\begin{array}{c|ccccc}
& e & z & C_1 & C_2 & C_3\\
\hline
\chi^{(1)} & 1 & 1 & 1 & 1 & 1\\
\chi^{(2)} & 1 & 1 & 1 & -1 & -1\\
\chi^{(3)} & 1 & 1 & -1 & 1 & -1\\
\chi^{(4)} & 1 & 1 & -1 & -1 & 1\\
\chi^{(5)} & 2 & -2 & 0 & 0 & 0
\end{array}
\]
(baris terakhirnya dari ortogonalitas kolom). Tabelnya identik untuk
$D_4$ dan $Q_8$, padahal keduanya tak isomorfik
(\cref{pb:b3:groups:1}). Tabel itu \emph{memang} menangkap: $\abs
G$, ukuran kelasnya, pusatnya ($\{g : \abs{\chi_i(g)} = n_i\
\forall i\}$: berorde $2$ pada keduanya), abelianisasinya, dan seluruh
kisi \omterm{def:b3:groups:normal}{subgrup normalnya} (lewat kernel dan irisannya,
\cref{exo:b3:representations:6}). Ia gagal menangkap orde
unsurnya: $D_4$ mempunyai lima involusi, sedangkan $Q_8$ satu --- jadi
tipe isomorfismanya sungguh lebih halus daripada \omterm{def:b3:representations:table}{tabel karakternya}.
\end{solution}

\begin{solution}{exo:b3:representations:5}
(a) Di $S_4$, \omterm{ex:b3:groups:actions}{sentralisator} $(1\,2\,3)$ berorde $24/8 =
3$: yakni $\langle(1\,2\,3)\rangle \subseteq A_4$. Jadi
$Z_{A_4}((1\,2\,3))$ berorde $3$ dan kelas-$A_4$ itu beranggota $12/3
= 4$ unsur: sehingga kedelapan siklus-$3$ terpecah menjadi dua kelas-$A_4$
(yang diwakili $(1\,2\,3)$ dan inversnya).

(b) Faktor $A_4/V \cong \Z/3\Z$ memberikan tiga \omterm{def:b3:representations:character}{karakter} berderajat $1$
(dengan $\omega = \eu^{2\iu\pi/3}$; kelas siklus-$3$ terpetakan ke
$\bar1, \bar2$); sedangkan pembatasan $\chi_{\mathrm{std}}$
tetap \omterm{def:b3:representations:rep}{tak tereduksi} ($\langle\chi,\chi\rangle = \frac1{12}(9 + 3
\cdot 1 + 0 + 0) = 1$):
\[
\begin{array}{c|cccc}
A_4 & e\,[1] & (1\,2)(3\,4)\,[3] & (1\,2\,3)\,[4] &
(1\,3\,2)\,[4]\\
\hline
\mathbf 1 & 1 & 1 & 1 & 1\\
\chi_\omega & 1 & 1 & \omega & \omega^2\\
\chi_{\bar\omega} & 1 & 1 & \omega^2 & \omega\\
\chi_3 & 3 & -1 & 0 & 0
\end{array}
\]
(c) Kernelnya: $\ker\chi_\omega = \ker\chi_{\bar\omega} = V$;
dan $\ker\chi_3 = \{e\}$ (sebab tak ada entri lain bermodulus $3$).
\omterm{def:b3:groups:normal}{Subgrup normalnya} berupa irisan kernelnya
(\cref{exo:b3:representations:6}): yakni $\{e\}$, $V$, $A_4$ ---
khususnya $A_4$ tak punya \omterm{def:b3:groups:normal}{subgrup normal} berorde $2$ atau berindeks
$2$.
\end{solution}

\begin{solution}{exo:b3:representations:6}
(a) Jika $\chi(g) = \chi(e) = n$: maka kesamaan pada
$\abs{\chi(g)} \leq n$ memaksa $\rho(g) = \lambda\,\mathrm{id}$
(\cref{prop:b3:representations:charbasics}) dengan $n\lambda = n$:
sehingga $\rho(g) = \mathrm{id}$. Kebalikannya jelas.

(b) Misalkan $N \trianglelefteq G$. \omterm{ex:b3:representations:examples}{Representasi reguler}
$G/N$ bersifat setia; uraikan ia atas \omterm{def:b3:representations:rep}{representasi tak tereduksi} $G/N$ lalu
angkat ke $G$ (\cref{prop:b3:representations:onedim}(b)):
diperoleh \omterm{def:b3:representations:character}{karakter} \omterm{def:b3:representations:rep}{tak tereduksi} $\chi_{i_1}, \dots$ dari $G$ yang kernelnya
memuat $N$ dan yang kernel \emph{persekutuannya} tepat berupa
prapeta $\{e\}$, yakni $N$ (karena kesetiaan pada \omterm{thm:b3:groups:quotient}{grup kuosiennya}).
Jadi $N = \bigcap_j \ker\chi_{i_j}$.

(c) Jika $G$ sederhana: untuk sebuah $\chi$ \omterm{def:b3:representations:rep}{tak tereduksi} yang tak trivial,
$\ker\chi \trianglelefteq G$ bukanlah $G$ (sebab \omterm{def:b3:representations:rep}{representasi tak
tereduksi} yang trivial pada seluruh $G$ adalah \omterm{def:b3:representations:character}{karakter} trivialnya),
sehingga $\ker\chi = \{e\}$. Sebaliknya, andaikan semua kernel tak
trivial itu trivial, dan misalkan $N \trianglelefteq G$ dengan $N \neq
G$. Pada ungkapan $N$ di (b) sebagai irisan kernel,
ada \omterm{def:b3:representations:character}{karakter} yang terlibat dan tak trivial (sebab kalau semuanya trivial,
irisannya akan $G$), dan kernelnya $\{e\}$: jadi $N =
\{e\}$. Jadi satu-satunya \omterm{def:b3:groups:normal}{subgrup normalnya} adalah $\{e\}$ dan $G$.
\end{solution}

\begin{solution}{exo:b3:representations:7}
Untuk orde $8$ yang tak abelian: banyaknya \omterm{def:b3:representations:character}{karakter} berderajat $1$ adalah
$[G : D(G)]$, yaitu pembagi sejati $8$ (sebab tak abelian: $D(G) \neq
\{e\}$), dan $\sum n_i^2 = 8$. Dengan $k$ buah angka satu dan sisa
derajatnya $\geq 2$: $8 - k \equiv 0$ dengan kuadrat $\geq 4$, dan $k
\mid 8$, $k < 8$. Untuk $k = 4$: satu derajat $2$ --- konsisten. Untuk $k =
2$: tersisa $6$, yang bukan jumlah kuadrat $\geq 4$. Untuk $k = 1$:
mustahil, sebab $k = [G : D(G)] \geq 2$ --- karena $G/D(G)$ merupakan
grup-$2$ abelian yang tak trivial, sebab $G$ grup-$2$ dengan
$D(G) \neq G$ menurut \omterm{def:b3:groups:derived}{keterselesaian} grup-$p$
(\cref{ex:b3:groups:solvableexamples}). Jadi polanya
$(1,1,1,1,2)$. Kesetiaan $\chi_5$: keempat \omterm{def:b3:representations:character}{karakter} berderajat $1$
semuanya memuat $D(G)$ di dalam kernelnya; seandainya $\ker\chi_5
\supseteq N \neq \{e\}$ untuk suatu $N$ normal minimal (entah $N
\subseteq D(G)$ entah tidak --- ambil $N \subseteq \ker\chi_5$
yang tak trivial), maka $N$ akan terletak di kelima kernelnya, yang
irisannya trivial (sebab \omterm{ex:b3:representations:examples}{representasi regulernya}
setia): kontradiksi. Untuk orde $p^3$ yang tak abelian: $Z(G)$
berorde $p$ (sebab orde $p^2$ akan membuat $G/Z$ siklik, sehingga $G$ abelian),
$G/Z(G)$ yang berorde $p^2$ bersifat abelian, sehingga $D(G) \subseteq Z(G)$,
dan $D(G) \ne \{e\}$: jadi $D(G) = Z(G)$, sehingga ada $p^2$ \omterm{def:b3:representations:character}{karakter}
berderajat $1$. Derajat sisanya memenuhi $\sum n_i^2 = p^3 -
p^2$ dengan setiap $n_i > 1$ membagi $\abs G$
(\cref{pb:b3:representations:1}, pertanyaan 8) sehingga $n_i \in
\{p\}$ (sebab $n_i = p^2$ akan melampaui: $p^4 > p^3 - p^2$): jadi tepat ada $p
- 1$ \omterm{def:b3:representations:character}{karakter} berderajat $p$.
\end{solution}

\begin{solution}{exo:b3:representations:8}
Untuk \omterm{def:b3:representations:rep}{representasi} $\rho, \sigma$ dari $G, H$ pada $V, W$, definisikan
\omterm{def:b3:representations:rep}{representasi} $\rho \boxtimes \sigma$ dari $G \times H$ pada
$V \otimes W$ --- secara konkret, pada matriksnya: $(\rho\boxtimes
\sigma)(g,h)$ adalah hasil kali Kronecker $\rho(g)\otimes
\sigma(h)$, yang tracenya
$\operatorname{tr}\rho(g)\operatorname{tr}\sigma(h) =
\chi(g)\psi(h)$ (sebab hasil kali Kronecker matriks $A \otimes B$
bertrace $\operatorname{tr}A\operatorname{tr}B$: diagonalnya
berupa $a_{kk}b_{ll}$). Jadi $\chi\psi$ merupakan \omterm{def:b3:representations:character}{karakter}, dan
\[
\langle\chi\psi, \chi'\psi'\rangle_{G\times H}
= \frac{1}{\abs G\abs H}\sum_{g,h}
\overline{\chi(g)\psi(h)}\,\chi'(g)\psi'(h)
= \langle\chi,\chi'\rangle_G\,\langle\psi,\psi'\rangle_H
= \delta_{\chi\chi'}\delta_{\psi\psi'}.
\]
Khususnya $\langle\chi\psi,\chi\psi\rangle = 1$: sehingga setiap
$\chi\psi$ \omterm{def:b3:representations:rep}{tak tereduksi}
(\cref{cor:b3:representations:multiplicity}). Inilah
$r_Gr_H$ \omterm{def:b3:representations:character}{karakter} \omterm{def:b3:representations:rep}{tak tereduksi} yang berbeda; sedangkan kelas
$G\times H$ berupa hasil kali kelasnya (sebab $(g,h) \sim (g',h')$
komponen demi komponen), sehingga ada $r_Gr_H$ kelas: jadi daftarnya
lengkap (\cref{thm:b3:representations:numberirr}). Untuk
$(\Z/2\Z)^2$: keempat \omterm{def:b3:representations:character}{karakter} tandanya
$(\pm1)\otimes(\pm1)$ --- yakni blok kiri atas tabel
\cref{exo:b3:representations:4}. Grup abelian
berhingga merupakan hasil kali grup siklik
(\cref{cor:b3:modules:abelian}); \omterm{def:b3:representations:character}{karakter} \omterm{def:b3:representations:rep}{tak tereduksinya} berupa
hasil kali \omterm{def:b3:representations:character}{karakter} sikliknya (\cref{exo:b3:representations:1}):
jadi semuanya berderajat $1$.
\end{solution}

\begin{solution}{exo:b3:representations:9}
(a) Berlaku $D(A_5) = A_5$ (sebab sederhana dan tak abelian), sehingga satu-satunya
\omterm{def:b3:representations:character}{karakter} berderajat $1$ adalah $\mathbf 1$
(\cref{prop:b3:representations:onedim}). Kita memerlukan $n_2^2 + n_3^2
+ n_4^2 + n_5^2 = 59$ dengan setiap $n_i \geq 2$; dengan menguji
kuadrat $4, 9, 16, 25, 36, 49$: satu-satunya multihimpunan yang berhasil adalah
$\{9, 9, 16, 25\}$: sebab dengan kuadrat terbesar $49$, sisanya $10$
bukan jumlah tiga kuadrat $\geq 4$; dengan $36$, sisanya
$23$ juga bukan ($16 + 4 + 4 = 24$, $9 + 9 + 4 = 22$); dengan
terbesar $25$, mudah diperiksa bahwa $25 + 16 + 9 + 9 = 59$ berhasil sedangkan semua ragam $25 +
25$, $25 + 16 + 16$, $25 + 16 + 4$ gagal; dan dengan terbesar
$16$: $16\cdot3 = 48 < 59 - 4$. Jadi derajatnya $1, 3, 3, 4, 5$.

(b) Permutasi pada $5$ titik: titik tetapnya $(5, 1, 2, 0, 0)$,
sehingga $\chi_4 = (4, 0, 1, -1, -1)$ dengan
$\langle\chi_4,\chi_4\rangle = \frac{16 + 0 + 20 + 12 + 12}{60}
= 1$: jadi \omterm{def:b3:representations:rep}{tak tereduksi}. Aksi pada keenam subgrup Sylow-$5$: sebuah
involusi menetapkan tepat $2$ di antaranya (sebab \omterm{ex:b3:groups:actions}{normalisatornya} dihedral berorde
$10$, yang masing-masing memuat $5$ involusi: jadi $30$ insidensi
untuk $15$ involusi), unsur berorde $3$ menetapkan $0$ (sebab tak ada orde $3$ di
$D_5$), dan unsur berorde $5$ menetapkan tepat $1$ (sebab ia terletak pada satu
Sylow saja): jadi \omterm{def:b3:representations:character}{karakter} permutasinya $(6, 2, 0, 1, 1)$, dan $\chi_5 =
(5, 1, -1, 0, 0)$ dengan norma $\frac{25 + 15 + 20 + 0 + 0}{60} =
1$: jadi \omterm{def:b3:representations:rep}{tak tereduksi}. Tersisa dua baris $(3, a, b, c, d)$, $(3, a', b', c',
d')$. Norma kolomnya
(\cref{cor:b3:representations:column}): pada kelas
$(1\,2)(3\,4)$ berlaku $\abs{Z} = 4$: $1 + 0 + 1 + a^2 + a'^2 = 4$;
kolomnya terhadap $e$: $1\cdot1 + 4\cdot 0 + 5\cdot1 + 3(a + a') =
0$: sehingga $a + a' = -2$, $a^2 + a'^2 = 2$: jadi $a = a' = -1$. Pada
siklus-$3$ berlaku $\abs Z = 3$: $1 + 1 + 1 + b^2 + b'^2 = 3$: jadi $b = b'
= 0$. Pada setiap kelas-$5$ berlaku $\abs Z = 5$: kolomnya terhadap $e$
terbaca $1\cdot 1 + 4\cdot(-1) + 5\cdot 0 + 3(c + c') = 0$, sehingga
$c + c' = 1$; dan
$c^2 + c'^2 = 3$: jadi $\{c, c'\} = \bigl\{\frac{1+\sqrt5}2,
\frac{1-\sqrt5}2\bigr\}$ --- rasio emas dan
sekawannya; sedangkan kelas-$5$ kedua membawa nilai yang tertukar (sebab
kedua barisnya harus ortogonal).

(c) Pada tabel yang sudah jadi, tak ada entri baris tak trivial yang sama dengan
derajatnya di luar kolom pertama: jadi setiap kernel $\{g :
\chi_i(g) = n_i\}$ bersifat trivial. Menurut
\cref{exo:b3:representations:6}(c), $A_5$ sederhana.
\end{solution}

\begin{solution}{exo:b3:representations:10}
Pemetaan $\rho(z)$ komutatif dengan setiap $\rho(g)$ (sebab $z$ di pusat), yakni
$\rho(z) \in \operatorname{End}_G(V) = \C\,\mathrm{id}$ (menurut Schur):
jadi $\rho(z) = \lambda_z\,\mathrm{id}$, dan $z \mapsto \lambda_z$ bersifat
multiplikatif: yakni morfisma $Z(G) \to \C^\times$. Lalu
$\chi(z) = n\lambda_z$ dengan $\abs{\lambda_z} = 1$ (karena akar
satuan): sehingga $\abs{\chi(z)} = n$. Jika $\rho$ setia, maka $\lambda$
injektif pada $Z(G)$ (sebab $\rho(z) = \mathrm{id} \iff \lambda_z =
1$), sehingga $Z(G)$ terbenam di $\C^\times$; dan subgrup berhingga dari
grup multiplikatif sebuah lapangan bersifat siklik
(\cref{thm:b3:galois:cyclic}).
\end{solution}

\begin{solution}{exo:b3:representations:11}
(a) Keekuivarianannya: $\rho(h)p_i\rho(h)^{-1}$ mengindeks ulang jumlahnya
(lewat $g \mapsto hgh^{-1}$, dan $\chi_i$ merupakan \omterm{def:b3:representations:character}{fungsi kelas}):
jadi $p_i$ komutatif dengan aksinya. Pada $W$ \omterm{def:b3:representations:rep}{tak tereduksi} yang
berkarakter $\chi_j$, Schur membuat pembatasannya menjadi skalar
$\lambda\,\mathrm{id}_W$; dengan mengambil tracenya,
\[
\lambda\,n_j = \frac{n_i}{\abs G}\sum_g
\overline{\chi_i(g)}\,\chi_j(g) = n_i\,\langle\chi_i,
\chi_j\rangle = n_i\,\delta_{ij}
\]
(lewat ortogonalitas pertama): jadi $\lambda = \delta_{ij}$.

(b) Uraikan $V$ atas \omterm{def:b3:representations:rep}{representasi tak tereduksi} (lewat Maschke): $p_i$ bekerja sebagai
identitas pada suku yang berkarakter $\chi_i$ dan sebagai $0$ pada
semua suku lain, sehingga $p_i$ merupakan proyeksi pada jumlahnya $V_i$
sepanjang jumlah sisanya; sedangkan peta $V_i$ tidak bergantung pada
penguraian yang dipilih (sebab ia himpunan vektor yang ditetapkan
$p_i$, terdefinisi tanpa pemilihan apa pun). Jumlah $\sum_ip_i$ bekerja sebagai identitas
pada setiap suku \omterm{def:b3:representations:rep}{tak tereduksinya}: jadi ia $\mathrm{id}_V$. Sedangkan
pemecahan yang lebih halus atas $V_i \cong W_i^{\oplus m_i}$ menuntut
pemilihan basis $\operatorname{Hom}_G(W_i, V)$: jadi ia tidak kanonik.

(c) Untuk $\varepsilon$ (berderajat $1$): $p_\varepsilon =
\frac1{6}\sum_{g}\varepsilon(g)\,\rho(g)$, yakni di dalam aljabar
grupnya
\[
p_\varepsilon = \tfrac16\bigl(e - (1\,2) - (1\,3) - (2\,3)
+ (1\,2\,3) + (1\,3\,2)\bigr) .
\]
Dengan mengkuadratkannya: koefisien $g$ pada $p_\varepsilon^2$ adalah
$\frac1{36}\sum_{h}\varepsilon(h)\varepsilon(h^{-1}g) =
\frac1{36}\,\varepsilon(g)\sum_h\varepsilon(h)^2 =
\frac{6}{36}\varepsilon(g)$: jadi $p_\varepsilon^2 =
p_\varepsilon$. (Petanya di dalam \omterm{ex:b3:representations:examples}{representasi reguler} berupa
garis yang direntang $\sum_g\varepsilon(g)e_g$: jadi \omterm{def:b3:representations:rep}{representasi}
tanda muncul dengan multiplisitas $1$, seperti dituntut teori
umumnya.)
\end{solution}

\begin{solution}{exo:b3:representations:12}
(a) Urutkan kelasnya $e$ [1], transposisi [6],
siklus-$3$ [8], siklus-$4$ [6], transposisi ganda [3].
\omterm{def:b3:representations:character}{Karakter} linear keduanya adalah $\chi_2 = \varepsilon$ dengan
nilai $1, -1, 1, -1, 1$. Untuk \omterm{def:b3:representations:character}{karakter} berderajat $2$, yakni
$\chi_3$, ortogonalitas kolom antara setiap kolom dengan
kolom identitasnya ($\sum_in_i\chi_i(g) = 0$ untuk $g \neq e$)
memberikan, pada transposisi: $1 - 1 + 2\chi_3 + 3(\chi_4 +
\chi_5) = 0$; lalu muslihat tanda $\chi_5 = \varepsilon\chi_4$
(sebab \omterm{def:b3:representations:character}{karakter} berderajat $3$ dikalikan \omterm{def:b3:representations:character}{karakter} linear tetap
\omterm{def:b3:representations:rep}{tak tereduksi} --- normanya sama) membuat $\chi_4 + \chi_5$ lenyap pada
kelas ganjil: jadi $\chi_3 = 0$ di sana. Pada siklus-$3$: $1 + 1 +
2\chi_3(c_3) + 3(\chi_4 + \chi_5)(c_3) = 0$ dengan
$\chi_5 = \chi_4$ pada kelas genap; sedangkan kolom $c_3$ dengan
dirinya sendiri memberikan data $\abs{\chi_3}^2$; dengan menyelesaikan sistem kecilnya
(pakai pula ortogonalitas baris $\chi_3$ dengan $\mathbf 1$ dan
$\varepsilon$): diperoleh $\chi_3 = (2, 0, -1, 0, 2)$, lalu $\chi_4 =
(3, 1, 0, -1, -1)$ dan $\chi_5 = \varepsilon\chi_4 = (3, -1,
0, 1, -1)$. Tabel selengkapnya:
\begin{center}
\begin{tabular}{c|ccccc}
 & $e$ & $6\,t$ & $8\,c_3$ & $6\,c_4$ & $3\,v$\\
\hline
$\chi_1$ & $1$ & $1$ & $1$ & $1$ & $1$\\
$\chi_2$ & $1$ & $-1$ & $1$ & $-1$ & $1$\\
$\chi_3$ & $2$ & $0$ & $-1$ & $0$ & $2$\\
$\chi_4$ & $3$ & $1$ & $0$ & $-1$ & $-1$\\
$\chi_5$ & $3$ & $-1$ & $0$ & $1$ & $-1$\\
\end{tabular}
\end{center}
(Semua barisnya bernorma $1$; semua kolomnya ortogonal:
jadi pemeriksaannya lulus.)

(b) Data kelas $1, 6, 8, 6, 3$ dengan derajat tersebut adalah data
milik $S_4$ (dengan tipe siklus $e$, $2$, $3$, $4$, $2{+}2$).
Kernelnya: $\ker\chi_2 = \{g : \varepsilon(g) = 1\} = A_4$
(kelas $e, c_3, v$: $1 + 8 + 3 = 12$, berindeks $2$);
dan $\ker\chi_3 = \{g : \chi_3(g) = 2\}$ = kelas $e, v$: yakni
grup Klein $V_4$, berorde $4$ dan normal. Sedangkan $\chi_4, \chi_5$
bersifat setia (sebab $\chi_i(g) = n_i$ hanya di $e$). Irisan
kernelnya: $\{e\}$, $V_4$, $A_4$, $G$ --- dan menurut
\cref{exo:b3:representations:6}(b) itulah \emph{semua}
\omterm{def:b3:groups:normal}{subgrup normal} $S_4$.

(c) \omterm{def:b3:representations:character}{Karakter} dengan $V_4 \subseteq \ker$ adalah $\chi_1,
\chi_2, \chi_3$: ketiganya terfaktorkan melalui $G/V_4$, yaitu grup
berorde $6$ yang berderajat \omterm{def:b3:representations:rep}{tak tereduksi} $1, 1, 2$ --- persis
tabel $S_3$. Karena tabel \omterm{thm:b3:groups:quotient}{grup kuosiennya} \emph{merupakan}
invarian lengkap di antara kedua grup berorde $6$
(sebab $\Z/6\Z$ akan mempunyai enam \omterm{def:b3:representations:character}{karakter} linear), maka $G/V_4 \cong
S_3$: jadi \omterm{thm:b3:groups:quotient}{grup kuosiennya} terlihat di dalam tabel itu sebagai blok
baris yang memuat $V_4$ di dalam kernelnya.
\end{solution}

\begin{solution}{pb:b3:representations:1}
\textbf{1.} Jika $\alpha^n + c_{n-1}\alpha^{n-1} + \dots + c_0 =
0$ (dengan $c_i \in \Z$), maka $\alpha^n \in \Z\text{-rentang}(1, \dots,
\alpha^{n-1})$, dan secara induktif setiap pangkatnya demikian: jadi $\Z[\alpha]$
dibangun oleh $1, \alpha, \dots, \alpha^{n-1}$. Sebaliknya misalkan
$\Z[\alpha] = \Z g_1 + \dots + \Z g_m$. Tulis $\alpha g_i =
\sum_j m_{ij}g_j$ dengan $M = (m_{ij}) \in M_m(\Z)$: maka vektor
$g = (g_i)$ memenuhi $(\alpha I - M)g = 0$; dengan mengalikannya dengan
matriks adjugatnya, $\det(\alpha I - M)\,g_i = 0$ untuk setiap $i$,
dan karena $1 \in \Z[\alpha]$ merupakan kombinasi-$\Z$ dari
$g_i$, maka $\det(\alpha I - M) = 0$: jadi $\alpha$ akar polinomial
monik $\det(XI - M) \in \Z[X]$.

\textbf{2.} Jika $\Z[\alpha]$ direntang oleh pangkat $\alpha$ sampai
$n-1$ dan $\Z[\beta]$ oleh pangkat $\beta$ sampai $m - 1$, maka
$\Z[\alpha,\beta]$ direntang oleh $nm$ hasil kali
$\alpha^i\beta^j$ (reduksikan sembarang monomialnya). Subring
$\Z[\alpha + \beta]$ dan $\Z[\alpha\beta]$ merupakan submodul-$\Z$
dari modul-$\Z$ $\Z[\alpha, \beta]$ yang \omterm{def:b3:modules:free}{dibangun secara berhingga},
sehingga keduanya pun \omterm{def:b3:modules:free}{dibangun secara berhingga} ($\Z[\alpha,\beta]$, yang dibangun oleh
$nm$ unsur, merupakan peta $\Z^{nm}$; sebuah submodul tertarik
kembali menjadi submodul $\Z^{nm}$, yang bebas berrank $\leq nm$ menurut
\cref{thm:b3:modules:submodule}, dan petanya membangunnya). Menurut
pertanyaan 1, $\alpha + \beta$ dan $\alpha\beta$ merupakan bilangan bulat
aljabar.

\textbf{3.} Jika $\frac pq$ (dalam bentuk paling sederhana) merupakan akar polinomial
bulat yang monik berderajat $n$, maka teorema akar rasional
(samakan penyebutnya: $p^n = -q(\cdots)$) memberikan $q \mid p^n$,
sehingga $q = \pm1$.

\textbf{4.} Nilai $\chi(g)$ merupakan jumlah akar satuan
(\cref{prop:b3:representations:charbasics}), yang masing-masing bilangan bulat
aljabar (sebab akar $X^N - 1$); lalu simpulkan dengan pertanyaan 2.

\textbf{5.} Untuk $h \in G$: $\rho(h)S_C\rho(h)^{-1} = \sum_{g\in
C}\rho(hgh^{-1}) = S_C$ (sebab $C$ sebuah kelas). Menurut Schur, $S_C =
\omega_C\,\mathrm{id}$; dan dengan mengambil tracenya, $\abs C\,\chi(g_C) =
\omega_C\, n$.

\textbf{6.} $S_CS_{C'} = \sum_{x \in C, y \in C'}\rho(xy) =
\sum_{z \in G} a(z)\rho(z)$ dengan $a(z) = \#\{(x,y) \in C\times
C' : xy = z\}$. Konjugasi dengan $h$ membijeksikan penyelesaian untuk
$z$ dengan penyelesaian untuk $hzh^{-1}$: jadi $a$ merupakan \omterm{def:b3:representations:character}{fungsi kelas} dengan
nilai di $\N$, sehingga $S_CS_{C'} = \sum_{C''}a_{CC'C''}S_{C''}$.
Dengan menyubstitusikan $S_C = \omega_C\,\mathrm{id}$ di semua tempat lalu
mengenali skalarnya:
$\omega_C\omega_{C'} = \sum_{C''}a_{CC'C''}\,\omega_{C''}$.

\textbf{7.} Misalkan $M$ modul-$\Z$ yang direntang oleh $1$ dan semua
hasil kali $\omega_{C_1}\cdots\omega_{C_k}$; menurut pertanyaan 6 setiap
hasil kali semacam itu tereduksi menjadi kombinasi-$\Z$ dari $1$ dan
$\omega_C$: jadi $M$ \omterm{def:b3:modules:free}{dibangun secara berhingga}, dan $\omega_C M
\subseteq M$ untuk setiap $C$. Khususnya $\Z[\omega_C]
\subseteq M$ \omterm{def:b3:modules:free}{dibangun secara berhingga} (sebagai submodul, seperti pada pertanyaan
2), lalu pertanyaan 1 menjadikan $\omega_C$ bilangan bulat aljabar.

\textbf{8.} Untuk sebuah $\chi$ \omterm{def:b3:representations:rep}{tak tereduksi} berderajat $n$:
\[
\frac{\abs G}{n} = \frac{\abs G}{n}\langle\chi,\chi\rangle
= \frac1n \sum_{g}\chi(g)\overline{\chi(g)}
= \sum_{C}\frac{\abs C\,\chi(g_C)}{n}\,\overline{\chi(g_C)}
= \sum_C \omega_C\,\overline{\chi(g_C)} .
\]
Setiap $\overline{\chi(g_C)} = \chi(g_C^{-1})$ merupakan bilangan bulat
aljabar (pertanyaan 4), sehingga ruas kanannya pun demikian (pertanyaan 2 dan
7); dan ia rasional, sehingga bilangan bulat (pertanyaan 3): jadi $n \mid \abs
G$.

\textbf{9.} Menurut Bézout: $u\abs C + vn = 1$ dengan $u, v \in \Z$.
Maka
\[
\frac{\chi(g_C)}{n} = u\,\frac{\abs C\,\chi(g_C)}{n} +
v\,\chi(g_C) = u\,\omega_C + v\,\chi(g_C),
\]
yang merupakan bilangan bulat aljabar.

\textbf{10.} Andaikan $0 < \abs{\chi(g_C)} < n$ lalu misalkan $\alpha
= \chi(g_C)/n$. Semua nilai $\chi(g)$ terletak di $\Q(\zeta_N)$ dengan $N =
\abs G$ (sebab berupa jumlah akar satuan ke-$N$). Untuk $\sigma \in
\operatorname{Gal}(\Q(\zeta_N)/\Q)$: $\sigma$ memetakan akar satuan
ke akar satuan (sebab $\sigma(\zeta^k) = \zeta^{ak}$,
\cref{thm:b3:galois:cyclotomicirred}), sehingga $\sigma(\chi(g_C))$
kembali menjadi jumlah $n$ akar satuan: jadi $\abs{\sigma(\alpha)}
\leq 1$; lagi pula $\sigma(\alpha)$ merupakan bilangan bulat aljabar (sebab ia
berpolinomial minimal sama dengan $\alpha$). Hasil kali $P =
\prod_\sigma \sigma(\alpha)$ ditetapkan oleh seluruh grup
Galoisnya, sehingga rasional (\cref{thm:b3:galois:fundamental}(1)),
dan ia bilangan bulat aljabar dengan
\[
0 < \abs P \leq \abs\alpha < 1
\]
(tak ada faktornya yang lenyap: sebab $\sigma(\alpha) = 0$ akan memaksa $\alpha =
0$). Ini bertentangan dengan pertanyaan 3. Karena itu $\chi(g_C) = 0$ atau
$\abs{\chi(g_C)} = n$, dan pada kasus terakhir $\rho(g_C)$ berupa
skalar (\cref{prop:b3:representations:charbasics}).

\textbf{11.} Ortogonalitas kolom (dengan $C \neq \{e\}$):
$\sum_i \chi_i(e)\overline{\chi_i(g_C)} = 0$, yakni $1 +
\sum_{i \geq 2} n_i\overline{\chi_i(g_C)} = 0$. Seandainya setiap
$\chi_i$ tak trivial dengan $p \nmid n_i$ lenyap di $g_C$, maka
dengan mengelompokkan sisanya menurut faktor $p$:
\[
-\frac1p = \sum_{i \geq 2,\ p \mid n_i}
\frac{n_i}{p}\,\overline{\chi_i(g_C)},
\]
akan menjadi bilangan bulat aljabar --- dan itu bertentangan dengan pertanyaan 3. Jadi ada
$\chi_i$ tak trivial dengan $p \nmid n_i$ dan $\chi_i(g_C) \neq 0$;
karena $\abs C = p^k$, berlaku $\gcd(\abs C, n_i) = 1$, sehingga pertanyaan 10
menjadikan $\rho_i(g_C)$ sebuah skalar. Sekarang $G$ sederhana dan tak abelian:
sehingga $\chi_i$ setia (\cref{exo:b3:representations:6}(c)), dan
$Z_i = \{g : \rho_i(g) \text{ skalar}\}$ merupakan \omterm{def:b3:groups:normal}{subgrup normal}
(yaitu prapeta skalarnya di bawah $\rho_i$, dan skalar membentuk
\omterm{def:b3:groups:normal}{subgrup normal} --- bahkan pusat --- dari petanya) yang memuat
$g_C \neq e$: jadi $Z_i = G$. Maka $\rho_i(G)$ abelian dan
setia, sehingga $G$ abelian: kontradiksi. \emph{Jadi tak ada \omterm{def:b3:groups:simple}{grup
sederhana} tak abelian yang mempunyai \omterm{ex:b3:groups:actions}{kelas konjugasi} berukuran pangkat prima
$> 1$.}

\textbf{12.} Misalkan $G$ sederhana berorde $p^aq^b$. Jika $b = 0$
(sehingga $G$ grup-$p$): maka $Z(G) \neq \{e\}$
(\cref{thm:b3:groups:pfixed}) bersifat normal, sehingga $Z(G) = G$: jadi abelian.
Jika tidak, ambillah $Q$ sebuah subgrup Sylow-$q$ dan $z \in Z(Q)
\setminus\{e\}$ (\cref{thm:b3:groups:pfixed} lagi). Maka $Q
\subseteq Z_G(z)$, sehingga kelas $z$ berukuran $[G : Z_G(z)]$
yang membagi $[G : Q] = p^a$: yakni pangkat $p$. Jika ukurannya $1$,
maka $z \in Z(G)$: pusatnya menjadi \omterm{def:b3:groups:normal}{subgrup normal} tak trivial, sehingga
$Z(G) = G$, jadi abelian. Jika ukurannya $p^k > 1$: pertanyaan 11
melarangnya bagi $G$ yang sederhana dan tak abelian. Bagaimanapun juga, \omterm{def:b3:groups:simple}{grup sederhana}
berorde $p^aq^b$ pastilah abelian (yakni $\cong \Z/p\Z$).

\textbf{13.} Induksi pada $\abs G$ (untuk $\abs G = 1$: \omterm{def:b3:groups:derived}{terselesaikan}). Jika
$G$ sederhana, pertanyaan 12 membuatnya abelian, sehingga \omterm{def:b3:groups:derived}{terselesaikan}.
Jika tidak, pilihlah $N$ \omterm{def:b3:groups:normal}{subgrup normal} sejati yang tak trivial: maka $\abs
N$ dan $\abs{G/N}$ kembali berbentuk $p^{a'}q^{b'}$ dan
lebih kecil, sehingga $N$ dan $G/N$ \omterm{def:b3:groups:derived}{terselesaikan} menurut induksi, dan $G$ pun
\omterm{def:b3:groups:derived}{terselesaikan} (\cref{prop:b3:groups:derived}). --- Dengan tiga
bilangan prima, langkah kuncinya gagal: sebab indeks sebuah \omterm{def:b3:groups:sylow}{subgrup Sylow} tak
lagi berupa pangkat prima, sehingga kelas unsur pusat sebuah
Sylow tak harus berukuran pangkat prima. Dan kegagalan itu memang
tak terelakkan: $A_5$, yang berorde $2^2\cdot3\cdot5$, sederhana dan
tak \omterm{def:b3:groups:derived}{terselesaikan}.

\textbf{14.} Kelas beserta ukurannya ada di
\cref{exo:b3:groups:11}(a). Kelas-$S_5$ dari $c$ berukuran
$24$; seandainya ia tetap satu kelas-$A_5$, pencacahan orbit--stabilisator
akan memberikan $\abs{Z_{A_5}(c)} = 60/24$, yang bukan bilangan bulat ---
secara konkret, $Z_{S_5}(c) = \langle c\rangle$ berorde $5$ dan
duduk di dalam $A_5$, sehingga kelas-$A_5$ dari $c$ berukuran $60/5 =
12$: jadi kelas-$S_5$ itu terpecah dua ($c$ dan $c^2$ sekawan
di $S_5$ hanya oleh permutasi ganjil).

\textbf{15.} Hanya ada satu \omterm{def:b3:representations:character}{karakter} linear: sebab \omterm{def:b3:representations:rep}{representasi} berderajat $1$
terfaktorkan melalui $G/D(G)$, sedangkan $D(A_5) = A_5$ (karena $A_5$ sederhana
dan tak abelian, dan $D(A_5)$ normal serta tak trivial ---
sebab $A_5$ tak abelian). Jadi $n_1 = 1$ dan $n_2^2 + n_3^2 + n_4^2
+ n_5^2 = 59$ dengan setiap $n_i \geq 2$. Kuadrat yang tersedia: $4, 9,
16, 25, 36, 49$. Untuk jumlah empat di antaranya yang sama dengan $59$:
terbesarnya haruslah $25$ ($36 + 4 + 4 + 9 = 53 < 59$ tak dapat
disesuaikan: $36 + 16 + 4 + 4 = 60$, $36 + 9 + 9 + 4 = 58$, $36 +
16 + 9 + 4 = 65$ --- tak ada kombinasi dengan $36$ atau $49$ yang berhasil),
dan $59 - 25 = 34 = 16 + 9 + 9$ (satu-satunya caranya: sebab $16 + 16 + 4 =
36$, $25 + 9 + 4 = 38$, $25 + 4 + 4 = 33$): jadi derajatnya $1, 3, 3,
4, 5$.

\textbf{16.} $\langle\pi, \mathbf 1\rangle = \frac1{60}(5 +
15\cdot1 + 20\cdot2 + 0 + 0) = 1$ (satu \omterm{def:b3:groups:action}{orbit} --- lewat cacahan
Burnside), dan $\langle\pi, \pi\rangle = \frac1{60}(25 + 15 + 80 +
0 + 0) = 2$: jadi $\pi$ memuat \omterm{def:b3:representations:character}{karakter} trivialnya sekali, dan
penyusun lainnya berupa satu \omterm{def:b3:representations:rep}{representasi tak tereduksi}. Karena itu $\chi_4 = \pi
- \mathbf 1$ \omterm{def:b3:representations:rep}{tak tereduksi}, berderajat $4$, dengan nilai $4, 0,
1, -1, -1$.

\textbf{17.} Pada pasangan, sebuah unsur menetapkan $\{i,j\}$ jika dan hanya jika ia menetapkan
atau menukar $i, j$: cacahannya $10$ (untuk $e$), $2$ (sebab $t =
(1\,2)(3\,4)$ menetapkan $\{1,2\}, \{3,4\}$), $1$ (sebab $(1\,2\,3)$
menetapkan $\{4,5\}$), lalu $0, 0$. Maka $\langle\pi_{10},
\pi_{10}\rangle = \frac1{60}(100 + 15\cdot4 + 20\cdot1) = 3$:
jadi tiga penyusun \omterm{def:b3:representations:rep}{tak tereduksi}, masing-masing sekali. Dan
$\langle\pi_{10}, \mathbf 1\rangle = \frac1{60}(10 + 30 + 20)
= 1$, $\langle\pi_{10}, \chi_4\rangle = \frac1{60}(40 + 0 +
20\cdot1\cdot1 + 0 + 0) = 1$: sehingga $\pi_{10} = \mathbf 1 + \chi_4
+ \chi$ dengan $\chi$ \omterm{def:b3:representations:rep}{tak tereduksi} berderajat $10 - 1 - 4 = 5$ dan
bernilai $\chi_5 = \pi_{10} - \mathbf 1 - \chi_4 = (5, 1, -1, 0,
0)$.

\textbf{18.} Tulis $a, a'$ untuk nilai $\chi_2, \chi_3$
di $t$, dan $b, b'$ di $s = (1\,2\,3)$; keempatnya real (sebab $t$
dan $s$ sekawan dengan inversnya). Kolom $t$ terhadap
kolom $e$: $1 + 3a + 3a' + 4\cdot0 + 5\cdot1 = 0$, sehingga $a + a'
= -2$; kolom $t$ dengan dirinya sendiri: $1 + a^2 + a'^2 + 0 + 1 =
\frac{60}{15} = 4$, sehingga $a^2 + a'^2 = 2$; karena itu $(a + a')^2 = 4
= 2 + 2aa'$ memberikan $aa' = 1$ dan $a = a' = -1$. Kolom $s$
terhadap $e$: $1 + 3(b + b') + 4 - 5 = 0$ memberikan $b + b' = 0$;
kolom $s$ dengan dirinya sendiri: $1 + b^2 + b'^2 + 1 + 1 = \frac{60}
{20} = 3$ memberikan $b = b' = 0$.

\textbf{19.} Kolom $c$ terhadap kolom $e$: $1 + 3(x + y) +
4(-1) + 5\cdot0 = 0$, sehingga $x + y = 1$. Kolom $c$ terhadap
kolom $c^2$ (kelas berbeda, jadi ortogonal): $1 + xy + yx + 1
+ 0 = 0$, sehingga $xy = -1$. Dengan demikian $x, y$ menyelesaikan $T^2 - T - 1 = 0$:
$\{x, y\} = \{\varphi, \bar\varphi\}$ dengan $\varphi = \frac{1 +
\sqrt5}2$. Konsistensinya: $x^2 + y^2 = (x+y)^2 - 2xy = 3 =
\frac{60}{12} - 2$ --- cocok dengan identitas kolom-diri $1 +
x^2 + y^2 + 1 + 0 = 5$. Tabelnya:
\begin{center}
\begin{tabular}{c|ccccc}
 & $e$ & $15\,t$ & $20\,s$ & $12\,c$ & $12\,c^2$\\
\hline
$\chi_1$ & $1$ & $1$ & $1$ & $1$ & $1$\\
$\chi_2$ & $3$ & $-1$ & $0$ & $\varphi$ & $\bar\varphi$\\
$\chi_3$ & $3$ & $-1$ & $0$ & $\bar\varphi$ & $\varphi$\\
$\chi_4$ & $4$ & $0$ & $1$ & $-1$ & $-1$\\
$\chi_5$ & $5$ & $1$ & $-1$ & $0$ & $0$\\
\end{tabular}
\end{center}

\textbf{20.} $\norm{\chi_2}^2 = \frac1{60}\bigl(9 + 15\cdot1 +
0 + 12\varphi^2 + 12\bar\varphi^2\bigr) = \frac{9 + 15 +
36}{60} = 1$, dengan memakai $\varphi^2 + \bar\varphi^2 = (\varphi +
\bar\varphi)^2 - 2\varphi\bar\varphi = 1 + 2 = 3$. Demikian pula
$\langle\chi_2, \chi_3\rangle = \frac1{60}(9 + 15 + 0 +
12(2\varphi\bar\varphi)) = \frac{9 + 15 - 24}{60} = 0$.
Keterbagian \omterm{thm:b3:galois:finitefields}{Frobenius}: $1, 3, 3, 4, 5$ semuanya membagi $60$.
Nilai irasional $\varphi, \bar\varphi$ itulah perbedaan yang
terlihat dengan $S_5$: sebab di $S_5$ setiap \omterm{def:b3:rings:divisibility}{unsur sekawan} dengan
semua pembangun grup sikliknya yang bertipe siklus sama
--- khususnya $c \sim c^2$ --- sehingga nilai \omterm{def:b3:representations:character}{karakternya} terpaksa rasional (bahkan
bulat); sedangkan di $A_5$ terpecahnya
siklus lima membuka pintu bagi $\Q(\sqrt5)$.

\textbf{21.} \omterm{def:b3:groups:normal}{Subgrup normal} $N$ berupa gabungan kelas, yang
memuat $e$, dengan $\abs N \mid 60$. Jumlah yang menjadi calonnya:
$1{+}15 = 16$, $1{+}20 = 21$, $1{+}12 = 13$, $1{+}24 = 25$,
$1{+}15{+}20 = 36$, $1{+}15{+}12 = 28$, $1{+}15{+}24 = 40$,
$1{+}20{+}12 = 33$, $1{+}20{+}24 = 45$, $1{+}12{+}12 = 25$,
$1{+}15{+}20{+}12 = 48$, $1{+}15{+}20{+}24 = 60 - 12 = \dots$
dengan mendaftarkan semua subjumlah sejati yang memuat $1$: tak satu pun dari $13, 16,
21, 25, 28, 33, 36, 40, 45, 48, 13{+}\dots$ yang membagi $60$
kecuali $1$ sendiri: jadi $N = \{e\}$ atau $A_5$. Kesederhanaannya terbaca dari
lima bilangan saja. Dan kriteria pertanyaan 11 juga terlihat: sebab
ukuran kelasnya $15 = 3\cdot5$, $20 = 4\cdot5$, $12 = 4\cdot3$ semuanya
komposit atas dua bilangan prima --- jadi tak ada kelas berukuran pangkat prima, persis
seperti dituntut kriteria Burnside bagi \omterm{def:b3:groups:simple}{grup sederhana}.

\textbf{22.} Rotasi $\R^3$ sebesar sudut $\theta$ mempunyai
nilai eigen $1, \eu^{\iu\theta}, \eu^{-\iu\theta}$: sehingga tracenya $1 +
2\cos\theta$. Untuk $\theta = \frac{2\pi}5$: $2\cos\frac{2\pi}5
= \frac{\sqrt5 - 1}2$, sehingga tracenya $1 + \frac{\sqrt5-1}2 =
\frac{1+\sqrt5}2 = \varphi$. Unsur tak trivial $\tau$ dari
$\operatorname{Gal}(\Q(\sqrt5)/\Q)$ yang diterapkan entri demi entri pada sebuah
\omterm{def:b3:representations:table}{tabel karakter} mengirim \omterm{def:b3:representations:character}{karakter} ke \omterm{def:b3:representations:character}{karakter} (sebab ia komutatif
dengan aljabar pendefinisinya: $\tau\circ\chi$ merupakan \omterm{def:b3:representations:character}{karakter}
\omterm{def:b3:representations:rep}{representasi} yang diperoleh dengan mengangkut matriksnya lewat
$\tau$ pada entrinya, atau secara abstrak: relasi ortogonalitasnya
rasional atas $\Q$, sehingga $\tau$ mempermutasikan penyelesaiannya);
lalu $\tau$ menetapkan $\chi_1, \chi_4, \chi_5$ (yang bernilai rasional) dan
karena itu harus menukar $\chi_2$ dan $\chi_3$: jadi \omterm{def:b3:representations:character}{karakter}
berderajat $3$ yang kedua membawa nilai sekawannya, tanpa satu matriks pun
dihitung. Secara geometri, kedua \omterm{def:b3:representations:rep}{representasi} itu adalah aksi
ikosahedral dan komposisinya dengan automorfisma luar
$A_5$ (yakni konjugasi dengan sebuah transposisi), yang
menukar kedua kelas siklus lima.

\textbf{23.} Untuk $z \in Z(G)$, $\rho(z)$ komutatif dengan setiap
$\rho(g)$, sehingga menurut lema Schur $\rho(z) = \lambda\,
\mathrm{id}$; dan karena $z$ berorde berhingga, $\lambda$ merupakan akar
satuan, sehingga $\abs{\chi(z)} = \abs\lambda\,n = n$. Lalu
\[
\abs G = \abs G\,\langle\chi, \chi\rangle
= \sum_{g \in G}\abs{\chi(g)}^2
\geq \sum_{z \in Z(G)}\abs{\chi(z)}^2
= \abs{Z(G)}\,n^2,
\]
yakni $n^2 \leq [G : Z(G)]$. Grup tak abelian berorde $8$
($D_4$ atau $Q_8$) mempunyai lima \omterm{ex:b3:groups:actions}{kelas konjugasi}, sehingga lima
derajat \omterm{def:b3:representations:rep}{tak tereduksi} dengan $\sum n_i^2 = 8$; dan satu-satunya cara
menulis $8$ sebagai jumlah lima kuadrat $\geq 1$ adalah $1 + 1 + 1 + 1
+ 4$: jadi derajatnya $1, 1, 1, 1, 2$. Kedua grup itu berpusat
berorde $2$, dan \omterm{def:b3:representations:character}{karakter} berderajat $2$ mencapai kesamaannya:
$2^2 = 4 = [G : Z(G)]$ --- jadi batasnya tajam. Untuk $A_5$,
$Z = \{e\}$ dan batasnya terbaca $n^2 \leq 60$: dipenuhi oleh
$1, 3, 3, 4, 5$ dengan sisa ruang yang lapang ($25 \leq 60$), dan memang harus
demikian sebab kesamaannya akan memaksa (lewat rantai yang sama) $\chi$
lenyap di luar pusatnya.

\textbf{24.} Berlaku $\rho(g)^m = \mathrm{id}$ untuk $m$ sama dengan orde
$g$, sehingga $\rho(g)$ dihapus oleh $X^m - 1$, yang terpecah dengan
akar sederhana atas $\C$: jadi dapat didiagonalkan, dengan nilai eigen
$\lambda_1, \dots, \lambda_n$ berupa akar satuan, pada sebuah
basis eigen $(e_i)$. Hasil kali $e_ie_j$ (dengan $i \leq j$) membentuk
basis eigen bagi $\operatorname{Sym}^2V$ dengan nilai eigen
$\lambda_i\lambda_j$, dan $e_i \wedge e_j$ (dengan $i < j$) membentuk basis eigen
$\Lambda^2V$; lalu karena
\[
\sum_{i<j}\lambda_i\lambda_j
= \frac{(\sum_i\lambda_i)^2 - \sum_i\lambda_i^2}2
= \frac{\chi(g)^2 - \chi(g^2)}2,
\]
sedangkan jumlah simetrisnya menambahkan $\sum_i\lambda_i^2$ alih-alih
menguranginya, maka kedua rumusnya menyusul. Untuk $\chi_2 = (3, -1,
0, \varphi, \bar\varphi)$ pada kelas $(e, (2,2)\text{-},
3\text{-siklus}, c, c^2)$: pengkuadratan mengirim transposisi
ganda ke $e$, siklus tiga ke siklus tiga, dan
kelas $c$ ke kelas $c^2$ serta sebaliknya (sebab $c^{-1} \sim
c$ di $A_5$, sehingga $c^4 \sim c$). Karena itu $\chi_2(g^2)$ terbaca $(3,
3, 0, \bar\varphi, \varphi)$, dan, dengan memakai $\varphi^2 = \varphi
+ 1$ serta $\bar\varphi = 1 - \varphi$:
\[
\Lambda^2\chi_2 = (3, -1, 0, \varphi, \bar\varphi) = \chi_2,
\qquad
\operatorname{Sym}^2\chi_2 = (6, 2, 0, 1, 1) .
\]
Dengan menguraikan yang terakhir itu, memakai ukuran kelas $1, 15, 20, 12,
12$: $\langle\cdot, \mathbf 1\rangle = \frac1{60}(6 +
15\cdot2 + 0 + 12 + 12) = 1$; $\langle\cdot, \chi_5\rangle =
\frac1{60}(30 + 30) = 1$; $\langle\cdot, \chi_4\rangle =
\frac1{60}(24 - 12 - 12) = 0$; $\langle\cdot, \chi_2\rangle =
\frac1{60}(18 - 30 + 12(\varphi + \bar\varphi)) = 0$, dan
demikian pula untuk $\chi_3$. Jadi $\operatorname{Sym}^2\chi_2 =
\mathbf 1 + \chi_5$ (dimensinya $6 = 1 + 5$) dan
$\chi_2\otimes\chi_2 = \mathbf 1 + \chi_2 + \chi_5$
(dimensinya $9 = 1 + 3 + 5$). Geometrinya: isomorfisma
ekuivarian $\Lambda^2\R^3 \to \R^3$, $u \wedge v \mapsto u
\times v$, tepat merupakan $\Lambda^2\chi_2 = \chi_2$ bagi sebuah
grup rotasi; sedangkan suku $\mathbf 1$ pada kuadrat
simetrisnya adalah bentuk kuadrat invarian $x^2 + y^2 + z^2$, dan
$\chi_5$ hidup pada ruang berdimensi lima berisi tensor
simetris bertrace nol (yaitu kuadrat harmonik).

\textbf{25.} Kolom $c$ terhadap kolom $c^2$:
\[
1\cdot1 + \varphi\bar\varphi + \bar\varphi\varphi +
(-1)(-1) + 0 = 1 - 1 - 1 + 1 + 0 = 0,
\]
seperti dituntut ortogonalitas bagi kelas yang berbeda (sebab $\varphi
\bar\varphi = -1$). Kolom $c$ terhadap dirinya sendiri: $1 +
\varphi^2 + \bar\varphi^2 + 1 + 0 = 1 + 3 + 1 = 5 = 60/12 =
\abs{Z_{A_5}(c)}$. \omterm{def:b3:representations:character}{Karakter} reguler $\sum_in_i\chi_i$ pada
keempat kolom bukan identitasnya:
\[
1 - 3 - 3 + 0 + 5 = 0, \qquad
1 + 0 + 0 + 4 - 5 = 0,
\]
pada transposisi ganda dan siklus tiga, lalu pada kelas
$c$ (sedangkan kelas $c^2$ merupakan sekawan Galoisnya):
\[
1 + 3\varphi + 3\bar\varphi - 4 + 0 = 1 + 3 - 4 = 0,
\]
dengan memakai $\varphi + \bar\varphi = 1$. Tabel itu lulus setiap
audit: jadi itulah \omterm{def:b3:representations:table}{tabel karakter} $A_5$.

\end{solution}
